\documentclass[11pt,a4paper]{article}

\usepackage[T1]{fontenc}
\usepackage[utf8]{inputenc}
\usepackage[brazilian, shorthands=off]{babel}
\usepackage{charter}
\usepackage[a4paper,margin=22mm]{geometry}
\usepackage{amsmath,amssymb}
\usepackage{booktabs}
\usepackage{array}
\usepackage{tabularx}
\usepackage{microtype}
\usepackage[hidelinks]{hyperref}

\setlength{\parskip}{6pt plus 2pt minus 1pt}
\setlength{\parindent}{0pt}

\newcommand{\code}[1]{{\ttfamily\footnotesize\allowbreak #1}}

\newcommand{\funcblock}[5]{%
  \vspace{5pt}\noindent{\raggedright\bfseries\ttfamily\footnotesize #1\par}%
  \vspace{1pt}\noindent\textit{#2}\par
  \begin{itemize}
    \setlength{\itemsep}{2pt}\setlength{\parskip}{0pt}
    \item \textbf{Parâmetros:} #3
    \item \textbf{Retorno:} #4
    \item \textbf{Exceções:} #5
  \end{itemize}
}

\sloppy

\begin{document}

% --- Página de Capa ---
\begin{titlepage}
  \centering
  \vspace*{1.2cm}

  {\scshape\large Fundação Getulio Vargas --- FGV EMAp\par}
  \vspace{0.25cm}
  {\normalsize Escola de Matemática Aplicada\par}
  \vspace{0.15cm}
  {\normalsize Análise Exploratória de Dados e Visualização\par}

  \vspace{3.0cm}

  \rule{\linewidth}{0.8pt}\par
  \vspace{0.6cm}
  {\LARGE\bfseries Documentação Técnica e Metodológica do Código\par}
  \vspace{0.45cm}
  {\large Dinâmica da Criminalidade no Rio de Janeiro:\\
  Uma Análise Integrada Sobre Segurança Pública\par}
  \vspace{0.6cm}
  \rule{\linewidth}{0.8pt}\par

  \vspace{2.6cm}

  {\normalsize\textbf{Autores}\par}
  \vspace{0.3cm}
  {\normalsize
    Josué Santos de Oliveira da Silva\\
    Leonardo Cardia Bicca\\
    Luigi Parmeggiani Porto\par
  }

  \vspace{1.2cm}

  {\normalsize\textbf{Professor}\par}
  \vspace{0.2cm}
  {\normalsize João Rulff\par}

  \vfill

  {\small\textbf{Repositório oficial:} [https://github.com/Josue-santoss/AEDV-A2.git]\par}
  \vspace{0.3cm}
  {\small Rio de Janeiro - 2026\par}
\end{titlepage}

\begin{abstract}
Esta documentação descreve a arquitetura computacional dos dois notebooks do projeto (\code{NotebookAEDV.ipynb} e \code{NotebookArtigo.ipynb}), o pipeline completo de preparação e tratamento de dados, as decisões metodológicas de código, de design visual e de interpretação descritiva, além da referência integral das 13 funções implementadas.
\end{abstract}

\tableofcontents
\vspace{8pt}

\section{Visão Geral, Arquitetura e Requisitos de Execução}

Todo o fluxo de ingestão, padronização, integração geoespacial, cálculo de indicadores e visualização interativa sobre os dados abertos do Instituto de Segurança Pública do Estado do Rio de Janeiro (ISP-RJ) é autossuficiente a partir do repositório oficial no GitHub ([https://github.com/Josue-santoss/AEDV-A2.git]), com suporte automático à leitura local da pasta \code{Projeto/Datas/} em execuções offline.

\textbf{Arquitetura dos artefatos do projeto:} o código está organizado em dois notebooks complementares que compartilham exatamente a mesma camada de dados, constantes e 13 funções auxiliares:
\begin{itemize}
  \setlength{\itemsep}{2pt}\setlength{\parskip}{0pt}
  \item \code{NotebookAEDV.ipynb}: notebook exploratório completo, estruturado em 15 etapas explícitas de preparação e auditoria de dados seguidas pelas sete perguntas de domínio, exibindo tabelas intermediárias, diagnósticos de cobertura e mapas interativos com mapa base (\code{px.choropleth\_map} sobre \code{carto-positron}).
  \item \code{NotebookArtigo.ipynb}: versão orientada à publicação em formato de artigo via \textbf{Quarto} (\code{execute: echo: false}, \code{embed-resources: true}), que consolida a preparação em três células iniciais, exibe tabelas-síntese selecionadas e emprega mapas vetoriais autocontidos (\code{px.choropleth} com projeção de Mercator e \code{pio.renderers.default = "notebook"}) para garantir renderização offline fiel em HTML.
\end{itemize}

\textbf{Ambiente e reprodutibilidade:} requer Python 3.10+ e as bibliotecas \code{requests}, \code{numpy}, \code{pandas}, \code{geopandas}, \code{matplotlib}, \code{seaborn}, \code{plotly}, \code{ipykernel} e \code{nbformat}. Para compilar \code{NotebookArtigo.ipynb} em HTML autocontido, requer adicionalmente o \textbf{Quarto CLI}.

\section{Preparação das Bases de Dados (Etapas 1 a 15)}

Esta seção documenta integralmente a \textbf{camada comum de dados} executada antes das perguntas de pesquisa e visualizações (detalhada em 15 etapas no \code{NotebookAEDV.ipynb} e condensada em três células no \code{NotebookArtigo.ipynb}). Cada tabela derivada possui uma granularidade definida (por exemplo, uma linha por \code{ano} e \code{id\_aisp}) e é reutilizada ao longo de todo o estudo.

\textbf{Princípios adotados:} preservar as bases originais em \code{dataframes\_brutos}, manter as versões tratadas em \code{dataframes}, separar limpeza de relacionamento territorial, construir indicadores uma única vez e registrar lacunas de cobertura sem inventar observações.

\subsection{Bibliotecas e configuração}
Centraliza as importações (\code{BytesIO}, \code{ZipFile}, \code{os}, \code{tempfile}, \code{unicodedata}, \code{warnings}, \code{requests}, \code{numpy}, \code{pandas}, \code{geopandas}, \code{matplotlib}, \code{seaborn} e \code{plotly}), a supressão do aviso \code{RuntimeWarning} referente aos arquivos de dados do GDAL, as constantes de acesso ao repositório oficial, a identidade visual acessível (paleta Okabe-Ito) e as cinco funções utilitárias básicas (\code{exigir\_colunas}, \code{taxa\_por\_100\_mil}, \code{agregar\_contagens\_aisp\_ano}, \code{formatar\_numero} e \code{adicionar\_fonte\_plotly}). O estilo global do Matplotlib remove bordas superior e direita, adota fundo branco e grade cinza clara; o Plotly utiliza o template \code{plotly\_white}. A Tabela~\ref{tab:constantes} resume as constantes globais definidas nessa etapa.

\begin{table}[!htbp]
\centering
\caption{Constantes globais de configuração e identidade visual.}
\label{tab:constantes}
{\footnotesize\begin{tabularx}{\linewidth}{>{\raggedright\arraybackslash}p{4.0cm} >{\raggedright\arraybackslash}p{4.4cm} >{\raggedright\arraybackslash}X}
\toprule
\textbf{Constante} & \textbf{Tipo / Valor} & \textbf{Finalidade} \\
\midrule
\code{REPOSITORIO\_DADOS} & \code{"Josue-santoss/AEDV-A2"} & Identificador \code{usuario/repositorio} no GitHub. \\
\code{BRANCH\_DADOS} & \code{"main"} & Branch de referência para CSVs e shapefiles. \\
\code{PASTA\_DADOS\_ZIP} & \code{"/Projeto/Datas/"} & Subcaminho filtrado dentro do ZIP baixado. \\
\code{PASTA\_DADOS\_LOCAL} & \code{"Projeto/Datas"} & Diretório local opcional utilizado em execuções offline. \\
\code{PREFIXO\_CSV} & \code{"br\_rj\_isp\_estatisticas\_"} \code{"seguranca\_"} & Prefixo removido ao gerar as chaves de \code{dataframes}. \\
\code{FATOR\_TAXA} & \code{100\_000} & Fator padrão de padronização por 100 mil habitantes. \\
\code{URL\_ZIP\_DADOS} & URL derivada & Endereço do arquivo ZIP do repositório. \\
\code{URL\_BASE\_SHAPEFILES} & URL derivada & Base \code{raw.githubusercontent.com} dos shapefiles. \\
\code{PALETA} & \code{dict[str, str]} & Paleta acessível Okabe-Ito e tons neutros. \\
\code{PALETA\_REGIOES} & \code{dict[str, str]} & Cores fixas para Capital, Baixada, Grande Niterói e Interior. \\
\code{FONTE\_PADRAO} & \code{str} & Atribuição padrão institucional (ISP-RJ). \\
\bottomrule
\end{tabularx}}
\end{table}

\subsection{Carregamento e catálogo das bases}
A função \code{carregar\_csvs\_do\_repositorio} lê prioritariamente os CSVs da pasta local \code{Projeto/Datas} (caso exista) ou baixa o arquivo \code{.zip} do GitHub em memória, identificando automaticamente todos os arquivos \code{.csv} dentro de \code{/Projeto/Datas/}, remove a extensão e o prefixo padrão e lê cada tabela com detecção automática de separador (\code{sep=None, engine="python"}), disparando \code{ValueError} em caso de nome duplicado ou ausência de arquivos.

São importadas \textbf{16 bases} em \code{dataframes\_brutos}: \code{PopulacaoEvolucaoMensalCisp}, \code{PopulacaoEvolucaoMensalMunic}, \code{armas\_apreendidas\_mensal}, \code{armas\_fogo\_apreendidas\_mensal}, \code{evolucao\_mensal\_cisp}, \code{evolucao\_mensal\_municipio}, \code{evolucao\_mensal\_uf}, \code{evolucao\_mensal\_upp}, \code{evolucao\_policial\_morto\_servico\_mensal}, \code{feminicidio\_mensal\_cisp}, \code{relacao\_cisp\_aisp\_risp}, \code{taxa\_evolucao\_anual\_municipio}, \code{taxa\_evolucao\_anual\_uf}, \code{taxa\_evolucao\_mensal\_municipio}, \code{taxa\_evolucao\_mensal\_uf} e \code{taxa\_letalidade}.

\subsection{Limpeza e padronização}
A rotina \code{limpar\_bases\_isp} gera \textbf{cópias} independentes em \code{dataframes}, mantendo \code{dataframes\_brutos} intacto. Para cada base:
\begin{enumerate}
  \setlength{\itemsep}{2pt}\setlength{\parskip}{0pt}
  \item Remove espaços nas extremidades dos nomes de colunas (\code{str.strip()}) e em colunas textuais (\code{object}/\code{string}) sem converter \code{NaN} na string \code{"nan"}.
  \item Corrige problemas de codificação, grafia e acentuação em \code{regiao} via \code{CORRECOES\_REGIAO} e normalização Unicode (\code{unicodedata.normalize("NFKD", ...)} com mapeamento canônico \code{REGIOES\_CANONICAS}), unificando variantes como \code{"Grande Niteroi"} em \code{"Grande Niterói"}.
  \item \textbf{Critério de imputação:} preenche ausências com zero \textbf{somente} em colunas numéricas iniciadas por \code{quantidade\_}, assumindo ausência de registro de ocorrência naquele recorte e preservando ausências em chaves, populações e taxas.
  \item Converte \code{ano} e \code{mes} para \code{int64} com \code{errors="raise"} e cria a coluna \code{data} (primeiro dia do mês) quando ambos existem.
\end{enumerate}

\subsection{Referências às bases de trabalho}
Define apelidos diretos para as nove tabelas efetivamente empregadas nas sete perguntas de domínio: \code{df\_armas}, \code{df\_cisp\_mensal}, \code{df\_mun\_mensal}, \code{df\_evol\_uf}, \code{df\_policial}, \code{df\_fem\_cisp}, \code{df\_rel}, \code{df\_taxa\_mun\_mensal} e \code{df\_letalidade}. Todas as 16 bases padronizadas permanecem disponíveis no dicionário \code{dataframes}, com inventário de dimensões (\code{linhas} e \code{colunas}) gerado em \code{inventario\_bases}.

\subsection{Dimensão de nomes das AISPs}
Centraliza uma fonte única de nomes de apresentação em \code{dim\_aisp} (colunas \code{id\_aisp}, \code{nome\_aisp} e \code{rotulo\_aisp} no formato \code{"AISP XX --- Nome"}), validando ausência de códigos duplicados. A tabela segue a referência oficial \textit{Relação RISP $\times$ AISP $\times$ CISP} do ISP-RJ (\url{https://www.ispdados.rj.gov.br/Arquivos/Relacao_RISPxAISPxCISP.pdf}) e inclui dois tratamentos históricos explícitos:
\begin{itemize}
  \setlength{\itemsep}{2pt}\setlength{\parskip}{0pt}
  \item \textbf{AISP 1:} identificada dinamicamente pelas CISPs observadas em cada ano, distinguindo a configuração histórica da Capital (CISPs 6 e 7 --- Cidade Nova / Santa Teresa, 2003--2011) e a reativação do 1º BPM em São Gonçalo (CISPs 72 e 73 --- São Gonçalo / Neves).
  \item \textbf{AISP 13:} cadastrada como \textit{Praça Tiradentes / Centro (extinta)}, referente ao extinto 13º BPM presente na série entre 2003 e 2011.
\end{itemize}

\subsubsection*{5.1. Chaves e nomes de colunas da população mensal}
Cria a cópia de trabalho \code{df\_populacao\_mensal\_cisp} a partir de \code{dataframes["PopulacaoEvolucaoMensalCisp"]}, valida a presença de \code{circ} e \code{pop\_circ} e espelha essas colunas nas chaves padronizadas \code{id\_cisp} e \code{populacao}.

\subsubsection*{5.2. Malhas geográficas (CISP, AISP, RISP)}
Carrega os shapefiles oficiais compactados (\code{CISP.zip}, \code{AISP.zip} e \code{RISP.zip}) por meio de \code{carregar\_shapefile\_github}, padronizando as chaves (\code{id\_cisp}, \code{id\_aisp}, \code{id\_risp}) como \code{int64}, reprojetando para \code{EPSG:4326} e simplificando as geometrias preservando topologia (\code{tolerance=0.005}), resultando em \code{gdf\_cisp} (137 polígonos), \code{gdf\_aisp} (41 polígonos) e \code{gdf\_risp} (7 polígonos).

\subsection{Relacionamentos territoriais (CISP, AISP e regiões)}
Constrói as correspondências territoriais \textbf{antes} de qualquer cálculo de taxas. O notebook valida \code{id\_cisp}, \code{id\_aisp} e \code{regiao} em \code{df\_rel} e lista as regiões cadastrais disponíveis em \code{aisps\_por\_regiao}.

\subsubsection*{6.1. Dimensão de relacionamento CISP $\rightarrow$ AISP}
Deriva \code{df\_mapa\_cisp\_aisp} do cadastro estático e verifica que nenhuma CISP está associada a mais de uma AISP. Para os indicadores históricos, constrói \code{mapa\_cisp\_aisp\_mensal} diretamente de \code{df\_cisp\_mensal} na granularidade \code{(ano, mes, id\_cisp)}, verificando ausência de chaves nulas e de dupla atribuição no mesmo mês. Constrói também \code{territorio\_aisp1}, registrando por ano as CISPs e os meses observados da AISP 1.

\subsubsection*{6.2. Região de policiamento predominante por AISP}
Enriquece \code{dim\_aisp} com a coluna \code{regiao\_predominante}, obtida pela moda da variável \code{regiao} em \code{df\_cisp\_mensal} para o período recente (\code{ano >= 2018}).

\subsection{Padronização da população mensal por CISP}
Na granularidade \code{(ano, mes, id\_cisp)}, substitui vírgula decimal por ponto em \code{df\_populacao\_mensal\_cisp["populacao"]}, converte para número com validação estrita contra valores ausentes ou negativos, trunca para \code{int64} e verifica a inexistência de registros duplicados para o mesmo \code{(ano, mes, id\_cisp)}.

\subsection{População anual por CISP}
No notebook, na granularidade \code{(ano, id\_cisp)}, calcula \code{df\_populacao\_anual\_cisp} como a \textbf{média} dos valores mensais disponíveis de população (sem somar populações entre meses).

\subsection{População anual por AISP}
Realiza junção \code{outer} (\code{validate="one\_to\_one"}) entre \code{mapa\_cisp\_aisp\_mensal} e \code{df\_populacao\_mensal\_cisp}, isolando eventuais registros sem vínculo territorial em \code{populacao\_cisp\_sem\_aisp}. Em cada mês, soma a população das CISPs da AISP; um mês só contribui para a média anual se todas as suas CISPs possuírem população (\code{cisps == cisps\_com\_pop}), recebendo \code{NaN} caso contrário. A tabela \code{populacao\_por\_ano\_aisp} guarda a média anual desses totais mensais completos, acompanhada de \code{meses\_ativos}, \code{meses\_com\_populacao} e \code{nome\_aisp} (com resolução temporal para a AISP 1). Cobertura populacional por AISP: \textbf{2003--2022}.

\subsection{População estadual}
Isola \code{populacao\_uf\_anual} a partir de \code{df\_letalidade} filtrando \code{regiao == "Estado"}, removendo duplicatas por \code{ano} e ordenando cronologicamente (cobertura: \textbf{1991--2022}).

\subsection{Contagens anuais de letalidade violenta por AISP}
Agrega \code{quantidade\_letalidade\_violenta} de \code{df\_cisp\_mensal} por \code{(ano, id\_aisp)} via \code{agregar\_contagens\_aisp\_ano}, produzindo \code{mortes\_aisp\_ano} (com mapeamento de \code{nome\_aisp} e tratamento temporal da AISP 1) e a série estadual agregada \code{letalidade\_anual}, sem restrição fixa de período.

\subsection{Painel AISP--ano e taxas por 100 mil habitantes}
Restringe \code{mortes\_aisp\_ano} aos anos com cobertura populacional e integra com \code{populacao\_por\_ano\_aisp} (\code{how="left", validate="one\_to\_one"}) formando \code{estatisticas\_aisp}. Calcula \code{taxa\_letalidade} com \code{taxa\_por\_100\_mil} exigindo \textbf{12 meses de atividade territorial} (\code{meses\_ativos == 12}) e população média positiva; anos com atividade parcial permanecem com \code{NaN} para evitar taxas artificiais.

\subsection{Verificações da cobertura do painel}
O \code{NotebookAEDV.ipynb} executa auditoria sistemática do relacionamento mortes $\times$ população: verifica a distribuição de \code{\_merge} (788 linhas \code{both}, 0 \code{left\_only}, 0 \code{right\_only}), lista as 8 linhas sem taxa por atividade parcial (por exemplo, AISP 1 em 2011 com 18 vítimas entre janeiro e junho e AISP 13 em 2003 e 2011), confirma 0 CISPs/meses sem vínculo territorial, valida a cobertura de nomes em \code{dim\_aisp} e monta \code{auditoria\_aisp1} acompanhando todos os anos com ocorrências da AISP 1. No \code{NotebookArtigo.ipynb}, a tabela \code{resumo\_cobertura} sintetiza a cobertura dos painéis integrados.

\subsection{Resumo anual das taxas por AISP}
Aplica \code{resumir\_taxas\_por\_ano} sobre \code{estatisticas\_aisp} para gerar \code{resumo\_estatistico}, contendo média, mediana, mínimo, máximo, quartis (Q1 e Q3), desvio-padrão populacional (\code{ddof=0}), variância, amplitude interquartil e coeficiente de variação percentual por ano. Essa tabela sintetiza a \textbf{dispersão entre AISPs} (heterogeneidade territorial), distinta da taxa estadual ponderada.

\subsection{Contrato de dados para as análises}
A Tabela~\ref{tab:contrato} consolida os objetos gerados na camada comum e suas granularidades para uso nas sete perguntas do projeto.

\begin{table}[!htbp]
\centering
\caption{Contrato de dados da camada comum reutilizada nas análises.}
\label{tab:contrato}
{\footnotesize\begin{tabularx}{\linewidth}{>{\raggedright\arraybackslash}p{4.3cm} >{\raggedright\arraybackslash}p{3.0cm} >{\raggedright\arraybackslash}X}
\toprule
\textbf{Objeto} & \textbf{Granularidade} & \textbf{Uso previsto} \\
\midrule
\code{dataframes\_brutos} & Depende da fonte & 16 CSVs originais para auditoria e reprocessamento. \\
\code{dataframes} & Depende da fonte & 16 bases padronizadas disponíveis para consulta. \\
\code{df\_rel} & Conforme arquivo ISP & Relações territoriais de origem. \\
\code{df\_mapa\_cisp\_aisp} & CISP & Cadastro de referência validado. \\
\code{mapa\_cisp\_aisp\_mensal} & Ano, mês e CISP & AISP observada no mês para indicadores históricos. \\
\code{dim\_aisp} & AISP & Nomes, \code{rotulo\_aisp} e \code{regiao\_predominante}. \\
\code{gdf\_cisp}, \code{gdf\_aisp}, \code{gdf\_risp} & Polígono territorial & Malhas geográficas padronizadas (\code{EPSG:4326}). \\
\code{df\_populacao\_mensal\_cisp} & Ano, mês e CISP & Série populacional mensal por delegacia. \\
\code{df\_populacao\_anual\_cisp} & Ano e CISP & Média anual das populações mensais por CISP. \\
\code{populacao\_por\_ano\_aisp} & Ano e AISP & Denominador anual para indicadores por AISP. \\
\code{populacao\_uf\_anual} & Ano & Denominador anual estadual (RJ). \\
\code{mortes\_aisp\_ano} & Ano e AISP & Contagem anual de letalidade violenta. \\
\code{estatisticas\_aisp} & Ano e AISP & Painel de métricas integradas com taxas válidas. \\
\code{resumo\_estatistico} & Ano & Distribuição anual das taxas entre AISPs. \\
\code{PALETA}, \code{PALETA\_REGIOES} & Constante global & Identidade visual acessível (Okabe-Ito). \\
\bottomrule
\end{tabularx}}
\end{table}

\section{Metodologia: Decisões de Código, Visualização e Interpretação}

Esta seção justifica as escolhas técnicas adotadas na implementação estatística/computacional, no desenho visual dos gráficos e nos limites interpretativos do estudo.

\subsection{Decisões de Engenharia de Código e Estatística}
\begin{itemize}
  \setlength{\itemsep}{3pt}\setlength{\parskip}{0pt}
  \item \textbf{Programação defensiva e integridade relacional:} todas as junções críticas especificam \code{validate="one\_to\_one"} ou \code{validate="many\_to\_one"}, e \code{exigir\_colunas} interrompe a execução diante de qualquer alteração de esquema nos CSVs, evitando multiplicação silenciosa de linhas.
  \item \textbf{Imputação seletiva e transparência de lacunas:} apenas contagens (\code{quantidade\_*}) recebem preenchimento com zero; denominadores populacionais nulos ou ausentes retornam \code{NaN}. Na Pergunta 4, a ausência de população estadual em 2023--2024 é tratada com \code{ffill()} (carregando a população de 2022) com aviso explícito em \code{NotebookAEDV.ipynb}.
  \item \textbf{Recorte por AISP, denominador feminino, pessoas-ano e expansão de registros (P1):} a opção por AISP compatibiliza a letalidade geral com a base de feminicídio (CISP/AISP) e desagrega a Capital. Na ausência de população por sexo, adotou-se a aproximação de 50\% de população feminina (\code{fator=200\_000}); para estabilizar eventos raros, os rankings usam taxas ponderadas por pessoas-ano em 2017--2022 exigindo \code{meses\_ativos == 12} em ambos os painéis (39 AISPs com cobertura completa entre os 41 polígonos da malha; AISPs 1 e 43 sem dados no período). Adicionalmente, a tabela diagnóstica \code{cobertura\_fem\_anual} contabiliza anualmente quantas AISPs registraram ao menos um feminicídio (2017--2025) para avaliar a hipótese de ampliação territorial da notificação.
  \item \textbf{Janela regional, dispersão e reconstrução de contagens (P2):} os nomes de região são normalizados via \code{normalizar\_regiao} (\code{unicodedata}) e a janela comum das quatro regiões é inferida automaticamente como \textbf{2000--2022} para o cálculo do coeficiente de variação regional (CV, que parte de 19,7\% em 2000, atinge pico de 36,5\% em 2014 e fecha em 22,5\% em 2022) e da participação percentual anual das ocorrências. Para evitar distorções de separador de milhar lido como decimal em \code{contagem\_delito}, as contagens regionais (\code{contagem\_p2}) são reconstruídas a partir de $\text{taxa} \times \text{população} / 100\,000$ com validação de soma igual a 100\%. Os extremos históricos são apurados sobre toda a série disponível (\textbf{1991--2022} para Capital e Estado; \textbf{2000--2022} para as demais), preservando \textbf{2003--2022} para variação acumulada e razão Máx/Mín.
  \item \textbf{Ajuste de calendário na sazonalidade (P3):} somas mensais brutas penalizam fevereiro (cerca de 9\% mais curto que meses de 31 dias). A conversão em média diária utilizando o número exato médio de dias de cada mês em 2003--2024 (incluindo os seis anos bissextos) e a normalização em Índice Base 100 removem esse viés de calendário; registra-se ainda que a soma acumulada pondera proporcionalmente mais os anos de maior volume histórico.
  \item \textbf{Clusterização robusta por medianas e teste de Spearman (P5):} a divisão em quadrantes pelas medianas de fuzis e de MDIP evita distorção por outliers; a correlação de Spearman sobre a taxa proporcional de fuzis confirma que a associação não é artefato do tamanho populacional.
  \item \textbf{Estimativa populacional municipal por identidade oficial (P6):} para estender a seleção dos 6 municípios mais populosos até 2024 (além do corte de 2022 das bases populacionais isoladas), inverte-se a identidade oficial entre contagem e taxa de registros de ocorrência.
  \item \textbf{Primeiras diferenças em defasagens temporais (P7):} para separar tendência secular compartilhada de resposta imediata no mês seguinte, calcula-se a correlação cruzada em primeiras diferenças mensais ($\Delta X_t$ vs.\ $\Delta Y_{t+k}$) contra a banda de referência exploratória aproximada sob hipótese de ruído branco ($\pm 1{,}96 / \sqrt{n_{\text{diff}}} \approx \pm 0{,}117$ para $n_{\text{diff}} = 283$, reportando o tamanho amostral efetivo por defasagem e ressalvando a ausência de pré-branqueamento formal).
\end{itemize}

A Tabela~\ref{tab:janelas} sintetiza as janelas temporais, granularidades e principais ressalvas metodológicas de cada uma das sete perguntas de domínio.

\begin{table}[!htbp]
\centering
\caption{Resumo de bases, granularidade, janelas temporais e ressalvas por pergunta.}
\label{tab:janelas}
{\footnotesize\begin{tabularx}{\linewidth}{>{\raggedright\arraybackslash}p{0.8cm} >{\raggedright\arraybackslash}p{3.0cm} >{\raggedright\arraybackslash}p{2.2cm} >{\raggedright\arraybackslash}p{2.3cm} >{\raggedright\arraybackslash}X}
\toprule
\textbf{Perg.} & \textbf{Base(s) principal(is)} & \textbf{Granularidade} & \textbf{Janela} & \textbf{Ressalva metodológica} \\
\midrule
P1 & \code{df\_cisp\_mensal}, \code{df\_fem\_cisp} & AISP--ano e UF--ano & 2003--2022 (let.); 2017--2022/25 (fem.) & Denominador feminino estimado em 50\%; 39 AISPs com série completa; ampliação gradual de registros. \\
P2 & \code{df\_letalidade} & Região--ano & 1991--2022 (comum 2000--2022) & Contagens reconstruídas de taxa $\times$ população; variação comparativa padronizada em 2003--2022. \\
P3 & \code{df\_cisp\_mensal} & Região--mês & 2003--2024 & Ajuste pela duração média de cada mês (Base 100); soma acumulada pondera mais anos de maior volume. \\
P4 & \code{df\_armas}, \code{df\_letalidade} & UF--ano & 2007--2024 & População estadual de 2022 propagada (\code{ffill}) para 2023--2024; apreensões refletem foco operacional. \\
P5 & \code{df\_cisp\_mensal}, \code{df\_armas} & AISP (acum.) & 2017--2022 & 39 AISPs com 12 meses ativos; quadrantes pelas medianas; associação territorial não implica causalidade. \\
P6 & \code{df\_mun\_mensal}, \code{df\_taxa\_mun\_mensal} & Município--ano & 2014--2024 & Top 6 municípios via inversão da taxa de ROs; MDIP compõe o denominador da letalidade violenta. \\
P7 & \code{df\_policial}, \code{df\_evol\_uf} & UF--mês (284m) & Jan/2003--Ago/2026 & Primeiras diferenças ($\pm 0{,}117$ sem pré-branqueamento); efeitos microlocais podem se diluir no estado. \\
\bottomrule
\end{tabularx}}
\end{table}

As formulações matemáticas centrais implementadas no código são:
\begin{align}
  T_{\text{letalidade}} &= \frac{\sum \text{ocorrências}}{\sum \text{população}} \times 100\,000, \qquad
  T_{\text{feminicídio}} = \frac{\sum F}{\sum (P / 2)} \times 100\,000 = \frac{\sum F}{\sum P} \times 200\,000, \\[4pt]
  I_{\text{sazonal}}(r, m) &= \frac{D_{r,m}}{\frac{1}{12}\sum_{j=1}^{12} D_{r,j}} \times 100, \quad \text{com } D_{r,m} = \frac{\text{Ocorrências}_{r,m}}{\text{Dias\_médios}_m}, \\[4pt]
  \widehat{P}_{\text{mun}} &= \frac{\text{ROs}}{\text{Taxa\_RO}} \times 100\,000, \qquad
  \text{Corr}_{\Delta}(k) = \text{Corr}(\Delta X_t,\, \Delta Y_{t+k}), \quad k \in \{0,1,2,3\}.
\end{align}

\subsection{Decisões de Design das Visualizações}
\begin{itemize}
  \setlength{\itemsep}{3pt}\setlength{\parskip}{0pt}
  \item \textbf{Acessibilidade e identidade cromática fixa:} adotou-se a paleta \textbf{Okabe-Ito} (segura para daltonismo) com mapeamento invariável por região (\code{PALETA\_REGIOES}) em todo o projeto.
  \item \textbf{Divisão entre bibliotecas estáticas e interativas:} mapas coropléticos (P1 e P5) e séries que se beneficiam de zoom, \textit{hover} unificado, animações ou filtragem por legenda (todos os gráficos da P2 --- trajetórias 1991--2022, coeficiente de variação 2000--2022 e participação regional animada por ano ---, P3, P6 e P7) foram construídos em \textbf{Plotly}; painéis comparativos densos (barras Top 10 e dispersões de P1, heatmap de P3, 2$\times$2 de P4, pequenos múltiplos de P6 e correlograma de P7) foram construídos em \textbf{Matplotlib/Seaborn}.
  \item \textbf{Implementação dos mapas entre os dois notebooks:} no \code{NotebookAEDV.ipynb}, os mapas de P1 e P5 utilizam \code{px.choropleth\_map} sobre o mapa base claro \code{carto-positron} (com o Cluster 4 em cinza médio \code{\#B0B0B0} para preservar contraste). No \code{NotebookArtigo.ipynb}, emprega-se \code{px.choropleth} com projeção de Mercator e \code{fitbounds="locations"}, dispensando servidores externos de \textit{tiles} para assegurar portabilidade no HTML gerado pelo Quarto. O mapa animado de letalidade em P1 cobre toda a janela populacional disponível (\textbf{2003--2022}), viabilizando a comparação direta de 2003, 2020 e 2022.
  \item \textbf{Evitar eixos Y duplos na P4 vs.\ escalas uniformes na P6:} na P4, volumes totais (milhares) e fuzis (centenas) foram separados em painéis próprios e comparados via \textbf{Índice Base 100}, evitando distorções de eixo duplo. Na P6, o uso de \code{twinx()} nos pequenos múltiplos 2$\times$3 emprega \textbf{escalas idênticas nos seis painéis} e rótulos explícitos em ambos os eixos verticais para preservar comparabilidade direta.
  \item \textbf{Controle de outliers, marcos históricos e acabamento:} nos mapas animados de P1, o teto da escala contínua (\code{YlOrRd} e \code{PuRd}) foi limitado ao \textbf{percentil 98} para preservar contraste territorial. Foram adicionados marcos históricos verticais (UPPs em 2008, Copa em 2014, Olimpíadas/crise fiscal em 2016 e Covid-19/ADPF 635 em 2020) em P2 e P4; contagem absoluta $n$ nas barras de feminicídio de P1 para sinalizar amostras pequenas; rotulagem seletiva apenas dos destaques do Cluster 1 no gráfico de dispersão de P5; remoção de bordas superiores/direitas e grades verticais; formatação numérica brasileira (\code{formatar\_numero}); e rodapé uniforme com a fonte ISP-RJ em todas as figuras (\code{adicionar\_fonte\_plotly}).
\end{itemize}

\subsection{Limites de Interpretação e Escopo Descritivo}
Em alinhamento com o princípio metodológico central do estudo, todos os indicadores e visualizações possuem caráter estritamente \textbf{descritivo e exploratório}:
\begin{itemize}
  \setlength{\itemsep}{2pt}\setlength{\parskip}{0pt}
  \item \textbf{Associação não implica causalidade:} quedas ou altas concomitantes (como o avanço das apreensões de fuzis durante a retração da letalidade violenta em P4, o pico do CV regional em 2014 em P2 ou a coexistência espacial no Cluster 1 em P5) subsidiam hipóteses descritivas --- como uma atuação policial mais incisiva contra o crime organizado ou o efeito de redistribuição territorial --- sem comprovar causalidade isolada.
  \item \textbf{Natureza administrativa dos registros:} os dados do ISP-RJ derivam de Registros de Ocorrência (ROs) e refletem conjuntamente a dinâmica criminal, práticas de notificação/tipificação (como a ampliação territorial dos registros de feminicídio após a Lei nº 13.104/2015) e o foco operacional das polícias.
  \item \textbf{Padrões sem causa estabelecida:} na sazonalidade (P3), a regularidade de picos em fevereiro--março e vales no inverno é documentada como padrão empírico recorrente sem atribuição causal conclusiva.
\end{itemize}

\section{Referência Completa das Funções}

Todas as 13 funções do projeto seguem docstrings no padrão PEP 257 (estilo Google).

\subsection{Utilitários de Validação, Cálculo e Formatação}

\funcblock{exigir\_colunas(base: pd.DataFrame, colunas: list[str], nome\_base: str) -> None}
  {Valida a presença de todas as colunas obrigatórias antes de qualquer agregação ou cruzamento.}
  {\code{base} (DataFrame alvo), \code{colunas} (nomes exigidos), \code{nome\_base} (rótulo para diagnóstico).}
  {\code{None}.}
  {\code{KeyError}, listando em ordem alfabética as colunas ausentes.}

\funcblock{taxa\_por\_100\_mil(quantidade: pd.Series, populacao: pd.Series, fator: int = FATOR\_TAXA) -> pd.Series}
  {Calcula taxas padronizadas por população, retornando \code{NaN} quando o denominador é ausente, nulo ou negativo.}
  {\code{quantidade} (contagens), \code{populacao} (denominador), \code{fator} (padrão \code{100\_000}; \code{200\_000} no feminicídio).}
  {\code{pd.Series} numérica com as taxas calculadas.}
  {Nenhuma (tratamento seguro via \code{where(populacao > 0)}).}

\funcblock{agregar\_contagens\_aisp\_ano(base: pd.DataFrame, colunas\_contagem: list[str]) -> pd.DataFrame}
  {Agrega colunas de contagem por \code{(ano, id\_aisp)} via soma após validar o esquema de entrada.}
  {\code{base} (DataFrame de ocorrências), \code{colunas\_contagem} (lista de métricas numéricas).}
  {\code{pd.DataFrame} com uma única linha por par \code{(ano, id\_aisp)}.}
  {\code{KeyError} via \code{exigir\_colunas}.}

\funcblock{formatar\_numero(valor: float, casas: int = 1, sinal: bool = False) -> str}
  {Formata números no padrão brasileiro (ponto para milhar e vírgula decimal) para rótulos gráficos Matplotlib.}
  {\code{valor} (número), \code{casas} (casas decimais), \code{sinal} (exibe \code{+} em positivos se \code{True}).}
  {\code{str} formatada.}
  {Nenhuma.}

\funcblock{adicionar\_fonte\_plotly(fig: go.Figure, texto: str = FONTE\_PADRAO, y: float = -0.16) -> go.Figure}
  {Insere anotação padronizada de fonte institucional no rodapé de figuras interativas Plotly.}
  {\code{fig} (figura alvo), \code{texto} (atribuição), \code{y} (posição vertical em coordenadas \code{paper}).}
  {A própria instância \code{go.Figure} atualizada.}
  {Nenhuma.}

\subsection{Ingestão, Limpeza e Sumarização}

\funcblock{carregar\_csvs\_do\_repositorio(url\_zip: str, pasta\_dados: str, prefixo: str, pasta\_local: str | None = PASTA\_DADOS\_LOCAL) -> dict[str, pd.DataFrame]}
  {Lê os CSVs prioritariamente da pasta local (se disponível) ou baixa o arquivo \code{.zip} do GitHub em memória com detecção automática de separador.}
  {\code{url\_zip} (URL do ZIP), \code{pasta\_dados} (subcaminho filtrado), \code{prefixo} (prefixo removido dos nomes), \code{pasta\_local} (diretório local opcional).}
  {\code{dict[str, pd.DataFrame]} mapeando nome padronizado da base ao DataFrame bruto.}
  {\code{requests.HTTPError}; \code{ValueError} se houver chave duplicada ou nenhum CSV.}

\funcblock{limpar\_bases\_isp(bases: dict[str, pd.DataFrame]) -> dict[str, pd.DataFrame]}
  {Produz cópias tratadas das bases: limpa espaços, corrige \code{regiao}, imputa zero em \code{quantidade\_*}, tipifica \code{ano}/\code{mes} e deriva \code{data}.}
  {\code{bases} (dicionário de DataFrames brutos).}
  {\code{dict[str, pd.DataFrame]} com todas as bases padronizadas.}
  {\code{ValueError} se \code{ano}, \code{mes} ou \code{data} contiverem valores inválidos.}

\funcblock{carregar\_shapefile\_github(nome\_zip: str, chave\_padrao: str, tolerancia: float = 0.005, pasta\_local: str | None = PASTA\_DADOS\_LOCAL) -> gpd.GeoDataFrame}
  {Lê shapefile compactado da pasta local (se disponível) ou baixa via arquivo temporário seguro, padroniza a chave territorial como \code{int64}, reprojeta para \code{EPSG:4326} e simplifica geometrias.}
  {\code{nome\_zip} (arquivo ZIP), \code{chave\_padrao} (ex.: \code{"id\_aisp"}), \code{tolerancia} (em graus; \code{None} desativa simplificação), \code{pasta\_local} (diretório local opcional).}
  {\code{gpd.GeoDataFrame} pronto para mapas coropléticos.}
  {\code{requests.HTTPError}, \code{KeyError} (chave ausente) e \code{ValueError} (CRS indefinido).}

\funcblock{resumir\_taxas\_por\_ano(painel: pd.DataFrame, coluna\_taxa: str = "taxa\_letalidade") -> pd.DataFrame}
  {Computa a distribuição anual entre AISPs: média, mediana, mínimo, máximo, Q1, Q3, desvio-padrão (\code{ddof=0}), variância, amplitude interquartil e coeficiente de variação (\%).}
  {\code{painel} (DataFrame AISP--ano), \code{coluna\_taxa} (métrica alvo).}
  {\code{pd.DataFrame} com uma linha por ano.}
  {\code{KeyError} via \code{exigir\_colunas}.}

\subsection{Funções Auxiliares de Visualização, Normalização e Classificação}

\funcblock{normalizar\_regiao(valor: object) -> str}
  {Normaliza o nome de uma região na Pergunta 2 removendo acentos (\code{unicodedata.normalize("NFKD", ...)}), padronizando caixa e removendo espaços extras.}
  {\code{valor} (texto ou objeto representando o nome da região na base de dados).}
  {\code{str} normalizada em minúsculas, sem acentuação e sem espaços nas extremidades.}
  {Nenhuma.}

\funcblock{anotar\_ponto(ax: plt.Axes, df: pd.DataFrame, col\_y: str, idx: int, texto: str, dy: float = 10, dx: float = 0, cor: str = PALETA["cinza\_escuro"], ha: str = "center") -> None}
  {Destaca picos e vales nos gráficos de linha da Pergunta 4 com caixa branca semitransparente e alinhamento horizontal configurável.}
  {\code{ax}, \code{df}, \code{col\_y}, \code{idx}, \code{texto}, \code{dy}, \code{dx}, \code{cor}, \code{ha}.}
  {\code{None}.}
  {Nenhuma.}

\funcblock{adicionar\_marcos\_historicos(ax: plt.Axes, marcos: list[tuple[int, str]], exibir\_rotulos: bool = False, y\_frac: float = 0.96, va: str = "top") -> None}
  {Desenha linhas verticais pontilhadas e rótulos contextuais (UPPs em 2008, crise fiscal em 2016 e Covid-19/ADPF 635 em 2020) nos quatro painéis da Pergunta 4.}
  {\code{ax}, \code{marcos} (lista de tuplas \code{(ano, rotulo)}), \code{exibir\_rotulos}, \code{y\_frac}, \code{va}.}
  {\code{None}.}
  {Nenhuma.}

\funcblock{classificar\_cluster(row: pd.Series) -> str}
  {Atribui cada AISP a um dos quatro quadrantes da Pergunta 5 comparando \code{quantidade\_fuzil} e \code{taxa\_mdip\_100k} às respectivas medianas.}
  {\code{row} (linha de \code{painel\_p5}).}
  {\code{str} com o rótulo do cluster (\code{Cluster 1} a \code{Cluster 4}).}
  {Nenhuma.}

\end{document}
